\documentclass[a4paper, 10pt]{article}
\usepackage[utf8]{inputenc}
\usepackage[pdftex]{graphicx}
\usepackage{tabularx}
\usepackage{amsmath}
\usepackage{amssymb}
\usepackage{geometry}
\usepackage{booktabs}
\usepackage{hyperref}
\usepackage{verbatim}
\usepackage{pdfpages}
\geometry{hmargin=2cm,vmargin=1.5cm}
%\renewcommand{\baselinestretch}{0.96}
\newcommand{\IR}{\mathbb{R}}
\newcommand{\IN}{\mathbb{N}}

\begin{document}

\begin{titlepage}
\begin{center}

{\Large Université de Mons}\\[1ex]
{\Large Faculté des Sciences}\\[1ex]
{\Large Département des mathématiques}\\[2.5cm]

\newcommand{\HRule}{\rule{\linewidth}{0.3mm}}
% Title
\HRule \\[0.3cm]
{ \LARGE \bfseries Rapport du TP2 \\[0.3cm]}
{ \LARGE \bfseries Analyse Numérique \\[0.1cm]}
\HRule \\[1.5cm]

% Author and supervisor
\begin{minipage}[t]{0.45\textwidth}
\begin{flushleft} \large
\emph{Professeur:}\\
Christophe \textsc{Troestler} \\
Quentin \textsc{Lambotte}
\end{flushleft}
\end{minipage}
\begin{minipage}[t]{0.45\textwidth}
\begin{flushright} \large
\emph{Auteurs:} \\
Loïc \textsc{Dupont} \\
Paolo \textsc{Marcelis} \\
Maximilien \textsc{Vanhaverbeke}
\end{flushright}
\end{minipage}\\[2ex]

\vfill

% Bottom of the page
\begin{center}
\begin{tabular}[t]{c c c}
\includegraphics[height=1.5cm]{logoumons.jpg} &
\hspace{0.3cm} &
\includegraphics[height=1.5cm]{logofs.jpg}
\end{tabular}
\end{center}~\\
 
{\large Année académique 2020-2021}

\end{center}
\end{titlepage}

\renewcommand{\contentsname}{Table des matières}
\tableofcontents

\newpage

\noindent
Soit \( M, N \in \IN \). Soit \( A \in \IR^{M \times N} \) avec rang\( A = N \leq M \). Soit \( b \in \mathbb{R}^n \). \\
Intéressons-nous à la résolution de l'équation \( Ax = b \). \\
Les exercices suivants consistent à regarder la meilleur solution au sens des moindres carrés càd la valeur de x qui minimise la fonction :
\begin{equation}
    \label{eq:1}
    x \mapsto |Ax - b|_2
\end{equation}
Pour les questions suivantes, nous allons utiliser les notations ci-dessous :
$$
A =
\begin{pmatrix}
a_{1, 1} & \cdots & a_{1, N} \\
\vdots   &        & \vdots \\
a_{M, 1} & \cdots & a_{M, N}
\end{pmatrix}
= (a_{i, j})_{
\tiny \begin{matrix}
1 \leq i \leq M \\
1 \leq j \leq N
\end{matrix}}
, \quad b =
\begin{pmatrix}
b_1 \\
\vdots \\
b_m
\end{pmatrix}
= (b_i)_{1 \leq i \leq M}, \quad x =
\begin{pmatrix} x_1 \\
\vdots \\
x_n
\end{pmatrix}
= (x_j)_{1 \leq j \leq N}
$$

\section{Exercice 1}

Soit \( x \in \IR^n \). On doit montrer que \( x \) réalise le minimum de \( x \mapsto |Ax - b|_2 \) \eqref{eq:1} si et seulement si \( x \) réalise le minimum de
\begin{equation}
    \label{eq:2}
    x \mapsto |Ax - b|_2^2
\end{equation}
Commençons par montrer que si \( x \) réalise le minimum de \eqref{eq:1}, alors \( x \) réalise le minimum de \eqref{eq:2}. \\
Supposons que \( x \) réalise le minimum de \eqref{eq:1} càd
$$
\forall y \in \IR^n,~ |Ax - b|_2 \leq |Ay - b|_2
$$
Comme toute norme est définie positive, on sait que pour tout \( z \in \IR^m,~ |z|_2 \in [0, +\infty[ \). \\
Alors, comme chaque membre de l'inégalité ci-dessus est positif, par croissance de la fonction \( x \mapsto x^2 \) sur \( [0, +\infty[ \), on a
$$
\forall y \in \IR^n,~ |Ax - b|_2^2 \leq |Ay - b|_2^2
$$
càd \( x \) réalise le minimum de \eqref{eq:2}. \\
Il reste à montrer que si \( x \) réalise le minimum de \eqref{eq:2}, alors \( x \) réalise le minimum de \eqref{eq:1}. \\
Supposons que \( x \) réalise le minimum de \eqref{eq:2} càd
$$
\forall y \in \IR^n,~ |Ax - b|_2^2 \leq |Ay - b|_2^2
$$
Comme tout carré d'un nombre réel est positif, par croissance de la fonction \( x \mapsto \sqrt{x} \) sur \( [0, +\infty[ \), on a
$$
\forall y \in \IR^n,~ \sqrt{|Ax - b|_2^2} \leq \sqrt{|Ay - b|_2^2}
$$
Comme les normes sont définies positives, \( \sqrt{|Ax - b|_2^2} = |Ax - b|_2 \) et \( \forall y \in \IR^n,~ \sqrt{|Ay - b|_2^2} = |Ay - b|_2 \). \\
Par conséquent, on a bien
$$
\forall y \in \IR^n,~ |Ax - b|_2 \leq |Ay - b|_2
$$
càd \( x \) réalise le minimum de \eqref{eq:1}. \\
On a donc bien montré que \( x \) réalise le minimum de \eqref{eq:1}, alors \( x \) réalise le minimum de \eqref{eq:2}.

\newpage

\section{Exercice 2}

Soit \( x \in \IR^n \). On doit montrer que \( x \) réalise le minimum de \( x \mapsto |Ax - b|_2 \) \eqref{eq:1} si et seulement si \( x \) vérifie l'équation
\begin{equation}
    \label{eq:3}
    A^T A x = A^T b
\end{equation}
Commençons par montrer que si \( x \) réalise le minimum de \eqref{eq:1}, alors \( x \) vérifie l'équation \eqref{eq:3}. \\
Supposons que \( x \) réalise le minimum de \eqref{eq:1} càd
$$
\forall y \in \IR^n,~ |Ax - b|_2 \leq |Ay - b|_2
$$
Par l'exercice 1, on sait que c'est équivalent à dire que
$$
\forall y \in \IR^n,~ |Ax - b|_2^2 \leq |Ay - b|_2^2
$$
Montrons que la fonction \( f : x \mapsto |Ax - b|_2^2 \) admet des dérivées partielles en tout point de \( \IR^n \). \\
Soit \( y = (y_1, \cdots, y_n) \in \IR^n \).

\end{document}